% !TeX root = ../papers/nc-p3.tex
\section{平行二重色の二つの曲面と異方向境界の貼り合わせ}
\label{sec:parallel-symmetric-surface}
\begin{theorem}\label{thm:parallel-symmetric-surface}
平行二重色の台94,95,96の全係数は、位置別文字変更の後、次の対称テンソル $S(p,q)$ とその四元群捻りの二曲面で尽くされる。
順列で同じになる成分を一度ずつ書くと
\[
 S_{0001}=p,\quad S_{0111}=S_{0122}=S_{0023}=S_{2333}=1,\quad S_{2223}=q
\]
であり、他は零とする。$(p,q)\in k^2$ で、両方零、片方零、両方非零が順に台94,95,96を与える。
両曲面の全点が標準Koszul左右AS正則で最小周期1、反対階数は順に6と12である。
各曲面内での無印の添字保持同型は、$(p,q)$ の交換と同時符号反転だけで尽くされる。
異方向二重色族と交わるのは捻りの対角線 $p=q$ だけで、その点は同族の射影対 $\{\infty,p\}$ と一致する。
\end{theorem}

\begin{proof}[全係数の被覆と全点の正則性]
$S$ の総色は1であり、第二文字を1だけ平行移動すると総色零の三つの代表台に正確に一致する。
台ごとに両切断の小行列式の二項式をSmith標準形で解くと、対角倍率を除いた解は順に次元0,1,2のトーラスが二つずつで、
残る小行列式は恒等的に零である。実際、定理\ref{thm:klein-parallel-flags}と同じ指数計算で、
語上の位置別倍率の指数行列の階数は13、非零の $p,q$ の指数列を加えた階数は順に13,14,15になる。
従って各表示のパラメータトーラスは一成分上に全射する。捻りとの相違は以下の反対階数で確かめられる。

$S$ の四文字に対応する線形形式を $X,Y,Z,T$ とする。隣接切断の四ブロックの階数は $2,2,1,1$ である。
対称関係空間中に平方が入る条件は、次の四つの二次式の共通零点である。
\[
 L=\langle X^2-pY^2-T^2,\ Z^2-Y^2-qT^2,\ YT,\ XT-YZ\rangle .
\]
$Y=0$ で $T\ne0$ なら第四式から $X=0$ となり第一式に反する。$T=0$ で $Y\ne0$ でも同様で、
$Y=T=0$ なら $X=Z=0$ が従う。従ってこの四式は基点自由で定理\ref{thm:symmetric-frobenius-chart}を適用できる。
斉次二コサイクル捻りは左右のKoszul性と双対完全性を保つ。捻り符号の一位置巡回による比は
移した一文字と総色1の間の交代双指標、すなわち位置別対角倍率である。従って両方が周期1である。
\end{proof}

\begin{proof}[捻りでの四文字の復元]
反対切断の色0,1の行列式はともに $-4$、色2,3はそれぞれ階数2である。
左核を行列とみなし列を $3,2,1,0$ の順にすると、その一般元は
\[
 \begin{pmatrix}z&0&0&0\\x&-t&0&0\\0&0&z&0\\0&0&y&t\end{pmatrix}.
\]
従って階数1の射影点は $x$ 軸と $y$ 軸の二点だけである。元の文字順でこれは両側の双対空間の文字1,3の二直線を復元する。
右核でも同じ二直線が得られる（次の始点を取ってもよい）。各位置で、その共通零空間 $E=\langle e_0,e_2\rangle$ と、
商空間の二つの文字直線が定まる。

この二段のフィルトレーションの随伴次数で、商文字が三つある成分を見る。各配置で
$0111$ と $2333$ の順列だけがあり、一つの偶位置を除くどの三位置でも同じ商文字を使う。
従って四商での文字交換は全て同じで、そこに伴う偶文字の二直線とその対応も決まる。
二つの文字対の同時交換を除けば、残る変換は対角倍率と、位置 $j$ での
$e_1\mapsto b_j e_1+u_j$, $e_3\mapsto d_j e_3+v_j$ ($u_j,v_j\in E$) である。
商文字を二つだけ残す成分は零である。二つの文字1を残した場合、他の二位置 $j,l$ には
$e_0\otimes u_l$ と $u_j\otimes e_0$ の非零倍率付き和だけが寄与する。
よって全 $u_j$ は $e_0$ 方向で、$e_0$ の位置別倍率 $a_j$ で割ったその係数は
全ての二位置で和が零となる。標数零より全て零である。
文字3を残した場合も同様である。ここで $2333$ の四つの順列の符号は、
唯一の文字2の位置だけで決まり、あらかじめその位置の文字2の倍率に吸収できる。よって $v_j=0$ である。

二文字対の交換は四元群の平行移動で、その捻り符号の変化は位置別倍率で実現する。
これを除くと、位置別対角倍率の比は $0111,0122,0023,2333$ の各語の全順列で同じでなければならない。
同じ語に含まれる二文字を任意の二位置で交換でき、文字間のグラフが連結なので、倍率は位置別定数を除いて四位置共通である。
初めの四つの非零成分を1に保つ文字倍率の指数行列を $M$、$0001,2223$ の指数行列を $N$ とすると
\[
 2NM^{-1}=\begin{pmatrix}1&-1&3&-1\\-1&3&-1&1\end{pmatrix}.
\]
従って $p,q$ の両倍率は共通の符号であり、負号も実現する
（$M$ による写像はトーラス間の全射で、二行とも整数格子に属さず、その和は $NM^{-1}$ で整数になる）。
零係数があってもこの計算は同じである。
\end{proof}

\begin{proof}[対称表示での一般の分解]
対称完全交叉の同型条件は $L$ の合同であり、同値にsocle四次テンソル $S$ の共通四文字変換による射影同値である。
$G$ を交代四次収縮 $D$ の勾配、$B=S_{01\mid23}G_{01\mid23}^{t}$ とする。
$z=pq$ と置くと $D=8(2z+3)$。部分trace $\operatorname{tr}_{1}(B^3)$ は偶文字0,2および奇文字1,3の各組で一つずつの対角値を持ち、
その差は $2048z(z-1)$ である。この作用素は共通倍率と相似を受けるだけで、$D=0$ の場合にも使える。
よって $z\ne0,1$ では二平面の無順序分解が定まる。
奇側三次の二形式は $Y^3,T^3$ であり、その張る空間の純立方はちょうどこの二直線である。
偶側三次は $pX^3+3XZ^2,3X^2Z+qZ^3$。ここに $(uX+vZ)^3$ が入るには
$uv\ne0$, $u^2=pv^2$, $v^2=qu^2$ が必要で、$z\ne1$ に反する。
従って二平面の順序も決まり、奇側の二立方に掛かる偶文字から全四文字が復元される。
残る群は前の指数計算による交換と同時符号だけである。

部分traceがスカラーとなる $z=0,1$ はそれぞれ $\operatorname{tr}(B^2)/D^2=8/3,88/25$ で区別される。
以下この二つの値での合同を直接調べる。
\end{proof}

\begin{proof}[対称表示の二つの例外軌跡]
$E_0=YT$ を指定して合同を取ると、既に同じ交換と同時符号だけが残る。
実際、その二因子と共通核への $L$ の制限の二平方 $X^2,Z^2$ が決まる。
制限の核 $J=\langle YT,XT-YZ\rangle$ により、$X$ は $T$ と、$Z$ は $Y$ と組になる。
従って交換を除くと持ち上げは $X'=aX+uY+vT$, $Z'=cZ+wY+hT$, $Y'=bY,T'=dT$ である。
$X'T'-Y'Z'\in J$ から $ad=bc$, $v=w=0$、第一・第二平方の $XY,ZT$ 係数から $u=h=0$。
二平方の残る係数が $a^2=d^2$, $c^2=b^2$ を強制し、二パラメータの倍率は共通の符号になる。

$E_0$ の像の候補は階数2の式で、その核への $L$ の制限が二次元、制限写像の核の鉛筆には可逆行列がある。
最後の条件は $XT-YZ$ から従う。一般の行列は
\[
 H=\begin{pmatrix}a&0&0&s\\0&-pa-b&-s&t\\0&-s&b&0\\s&t&0&-a-qb\end{pmatrix}.
\]
$a=b=0$ なら階数2の候補は $E_0$ だけである。$a=0,b=1$ なら三次小行列式から $s=0,t^2=q$ となる。
その核への制限は左上が $a$、非対角が $s$、右下が $-(1+pt_0^2)a-2t_0^2 b+2t_0 t$ の二行列である
（$t_0$ は候補点の $t$）。$t_0\ne0$ なら制限は三次元、零なら $a=s=0$ の鉛筆全体が特異で、どちらも除かれる。
$b=0$ の場合も交換で同じである。

$ab\ne0$ ではSchur補行列から $t=0$ と
\[
 s^2=-b(pa+b)=-a(a+qb)
\]
が必要十分である。$a=1$ と置く。$s\ne0$ のとき $b^2+(p-q)b-1$ の根が単純なら、
二つの対角条件の $b,s$ に関する微分は独立である。さらに $t$ の微分を加えると核への制限は三次元である
（階数2点でSchur補行列の微分は核への制限そのもの）。$z=0$ で重根が出るのは $(0,\pm2\sqrt{-1})$ とその交換だけで、
これらはすでに同じ類である。$z=0$ の他の点には $s=0$ の候補もない。
指定した $E_0$ を保存しない同型があればその逆でも保存しないため、他の点からこの例外類への追加同一視もない。

$z=1$ で二次式が重根なら $p+1/p=0$ で $b=-p,s=0$。従って全ての $z=1$ で
$E_0$ 以外の候補は $X^2-pZ^2$ だけである。
その二因子は $X\pm\sqrt p Z$、核上の平方直線は $(\sqrt p Y\pm T)^2$ である。
一方、その制限の核を張るもう一式 $XT-YZ$ は、二因子 $X\pm\sqrt p Z$ に
$T\mp Y/\sqrt p$ を掛け合わせて得る。元の $E_0$ のところでは、
こうして得る二直線が制限の平方直線と一致する。従って新候補に移るには
$\{\pm\sqrt p\}=\{\pm1/\sqrt p\}$、つまり $p=\pm1$ が必要である。
逆同型にも同じ条件が必要で、その二点は $q=p$ で同時符号により同型である。これで全例外を尽くした。
\end{proof}

\begin{proof}[枝間の貼り合わせ]
対称表示は反対階数6で異方向族から分離される。捻りの左反対核の行列式は符号を除き $z^2t^2$ であるため、
定理\ref{thm:transverse-doublet-isomorphism}の二次錐対を持つ階数12枝とは交わらない。
残る $\{\infty,h\}$ の標準形は、全位置の文字に1を加えると、捻りを戻した表示でちょうど $S(h,h)$ になる。
平行移動の捻り符号の変化は位置別倍率なので、証明済みの無印同型により対角線だけが一致する。
従って二つの曲面はいずれも既約な二次元像を持ち、台95と94は各々の座標軸とその交点の特殊化である。
\path{computations/verify_parallel_symmetric_surface.py} が二成分の指数、切断、部分trace、例外の二次式、境界のテンソル一致を検算する。
\end{proof}
